%!TEX program = lualatex
\documentclass[ngerman,headheight=70pt]{scrartcl}
\usepackage[ngerman]{babel}
\usepackage{blindtext}
\usepackage[utf8]{luainputenc}
\usepackage{scrlayer-scrpage}
\usepackage[T1]{fontenc}
\usepackage{fontspec}
\usepackage{multicol}
\usepackage{enumitem}
\usepackage{epigraph}
\usepackage[german=quotes]{csquotes}
\usepackage[bottom=2cm,footskip=8mm]{geometry}
\setmainfont{Source Sans Pro}
\MakeOuterQuote{"}

\setlength{\parindent}{0pt}
\setlength{\parskip}{6pt}

\pagestyle{plain.scrheadings}
\thispagestyle{scrheadings}
\lohead{\vspace*{-1.5cm}\includegraphics[scale=0.4]{../images/stupa_logo_rot.pdf}}
\rohead{\\Präsidium des\\Studierendenparlaments}

\lofoot[Universität Hamburg · Präsidium des Studierendenparlaments\\\hrule Von-Melle-Park 5 · D-20146 Hamburg]{Universität Hamburg · Präsidium des Studierendenparlaments\\\hrule Von-Melle-Park 5 · D-20146 Hamburg}
\cofoot[]{}
\cohead[\thepage]{}

\renewcommand{\thesubsection}{TOP \arabic{subsection}}
\renewcommand{\thesubsubsection}{\arabic{subsubsection}.}
\setcounter{subsection}{-1}

\begin{document}
    UHH · StuPa-Präsidium · Von-Melle-Park 5 · D-20146 Hamburg

    \section*{Beschluss des Studierendenparlamentes vom 12. Mai 2016}
    \subsection*{Auseinandersetzung mit der AfD}

    Das StuPa bekräftigt, dass der Ausschuss gegen Rechts in Zusammenarbeit
    mit dem AStA einen "Tag gegen Rechts" entwickeln, planen und durchführen
    möge und wird dies nach Kräften unterstützen.

    Insbesondere möge vorbereitet werden, dass sich an diesem Tag mit der AfD
    kritisch-analytisch auseinandergesetzt wird.

    Auf der Grundlage der Erkenntnisse des "Tages gegen Rechts" soll der Ausschuss
    gegen Rechts eine politische Stellungnahme gegen die AfD ausarbeiten, die
    dem StuPa zur Beschlussfassung vorgelegt werden.
\end{document}
